\documentclass[12pt]{exam}
\usepackage[T1]{fontenc}
\usepackage[utf8]{inputenc}
\usepackage[lithuanian]{babel}
\usepackage{lmodern}
\usepackage{graphicx}
\usepackage[a4paper, left=1cm, right=1cm, top=2.5cm, bottom=1cm]{geometry}
\usepackage{amsmath, amssymb}
\usepackage{tasks}

% Antraštė
\pagestyle{headandfoot}
\runningheadrule
\firstpageheader{Vardas: \rule{4cm}{0.4pt}}{Kontrolinis darbas: \\Daugianariai}{Platono variantas}
\runningheader{Daugianariai}{}{}
\firstpagefooter{}{}{}
\runningfooter{}{}{}
\pointsdroppedatright          
\bracketedpoints               
\setlength{\rightpointsmargin}{1.6cm} 
\renewcommand{\points}{taškai} 
\renewcommand{\pointname}{Taškai:} 
\marginpointname{\ taškai}     
\renewcommand{\dottedlinefillheight}{0.7cm}
\newcommand{\atsakymas}{%
    \par\vspace{0.3cm}
    \textbf{Atsakymas:} \framebox[7cm]{\rule{0pt}{0.7cm}}
}

\begin{document}

\begin{questions}

\question[3] Nustatykite duotų reiškinių tipus (racionalusis, trupmeninis, sveikas, iracionalusis). Atsakymą pateikite prie kiekvieno reiškinio prirašydami raides pvz. R,S (Racionalus ir Sveikas).
\begin{tasks}[label-width=1.2em](3)
    \task $\displaystyle \frac{x^4 - \sqrt{5}}{3}$
    \task $\displaystyle \frac{4}{x^3 + 2}$
    \task $\displaystyle x^{-4} + \sqrt{2} x$
    \task $\displaystyle \sqrt{x^2 - 9}$
    \task $\displaystyle 2x+3 ^{-2} $
    \task $\displaystyle \sqrt{2}x-7 $
\end{tasks}
\droppoints

\question[4] Suprastinkite reiškinį ir jeigu atsakymas yra vienanaris, nurodykite jo koeficientą ir laipsnį. Jei ne vienanaris, paaiškinkite kodėl.
$$ \left( -2x^2y^{-3} \right)^5 \cdot \left( 2x^{-3}y^4 \right)^{-4} $$
\droppoints
\vspace{2.5cm}
\atsakymas

\question[4] Nenaudodami prietaiso, kuris homo sapiens paverčia homo babuinus, apskaičiuokite šio reiškinio reikšmę:
 $$\displaystyle \frac{68^3 + 32^3}{0,2 ^{-2} \cdot (68 \cdot 36+32 ^{2} ) } $$ \droppoints
\vspace{3.5cm}
\atsakymas

\question Įrodykite formules:
\begin{parts}
    \part[2] $\displaystyle \frac{x^m}{x^n} = x^{m-n} $ \droppoints
    \vspace{2cm}
    \part[3] $(x-y) ^{3} = x^3 - 3x^2y + 3xy^2 - y^3 $. \droppoints
\end{parts}
\newpage

\question[4] Raskite didžiausią ir mažiausią reiškinio $-5x^2 + 2 x + 1$ reikšmes. Nurodykite, su kokia $x$ reikšmėmis jos gaunamos.
\droppoints
\vspace{3.5cm}
\atsakymas

\question[5]  Reiškinį $(x+1) ^{3} -(x+2) ^{2} (x+1)$ užrašę pavidalu $a(x+b)^{2} + c$ apskaičiuote reiškinio $\sqrt{-2a-4b+8c+7} $ reikšmę.   
\droppoints
\vspace{3.5cm}
\atsakymas

\question[4] Kokia yra mažiausia galima reiškinio $x ^{2} +y ^{2} + z ^{2} +6x-4y-8z$ reikšmė? 
\droppoints
\vspace{1cm}

\end{questions}

\vfill 

\begin{center}
    \textbf{Vertinimo lentelė (Iš viso: 29 taškai)} \\[0.2cm]
    \renewcommand{\arraystretch}{1.3} 
    \begin{tabular}{|c|c|c|c|c|c|c|c|c|c|c|}
        \hline
        \textbf{Taškai} & 0--3 & 4--6 & 7--9 & 10--12 & 13--15 & 16--18 & 19--21 & 22--24 & 25--27 & 28--29 \\ \hline
        \textbf{Pažymys} & \textbf{1} & \textbf{2} & \textbf{3} & \textbf{4} & \textbf{5} & \textbf{6} & \textbf{7} & \textbf{8} & \textbf{9} & \textbf{10} \\ \hline
    \end{tabular}
\end{center}

\end{document}
